\section{Introduction}
\label{sec:introduction}
\IEEEPARstart{D}{ata} centres are large electricity consumers whose internal operating choices can also provide flexibility to the power system. Their computing equipment, backup batteries and cooling infrastructure offer different ways to change demand, but these resources serve a common purpose: maintaining the facility's computational service. Assessing their flexibility therefore requires more than identifying spare electrical capacity. It requires establishing how much the facility can change its grid demand, for how long, and what must happen afterwards to restore its operating condition.

The significance of this question is increasing with the scale and concentration of digital infrastructure. In its 2025 \emph{Energy and AI} report, the International Energy Agency estimated that data centres consumed 415 TWh in 2024 and projected approximately 945 TWh by 2030 in its base case~\cite{iea_energy_ai}. Alongside this demand growth, increasing renewable generation strengthens the need to adjust electricity consumption to changing supply conditions. The geographical concentration of data centres makes their local demand profiles important as well as their aggregate energy use. Flexible operation could help accommodate these loads, provided that the response available from a facility can be described in terms useful to a system operator.

Data-centre flexibility has both an operational and an economic dimension. Work can be moved between permitted execution times, batteries can exchange energy with the internal electrical system, and cooling can use stored thermal energy or limited temperature variation. Such actions may reduce electricity expenditure or change demand in response to a grid request. However, a schedule that is attractive under electricity prices does not necessarily leave every resource ready to deliver an additional response. A charged battery, available computing capacity and a favourable thermal state cannot be assumed to coincide. The operating trajectory from which a request begins is therefore part of the definition of available flexibility.

Throughout this paper, \emph{upward flexibility} means a reduction in electricity import, supporting the system when demand must fall relative to supply. \emph{Downward flexibility} means an increase in import. Requested deviations are signed from the consumer's perspective: negative for reduced import and positive for increased import. The study characterises their feasible duration and the contributions of internal resources. It does not assume that the facility exports electricity or that a response lasting several hours establishes eligibility for a particular frequency-response product.

\subsection{Data-centre architecture}
\label{sec:architecture}
A data centre combines IT equipment with the electrical and thermal infrastructure needed to sustain it. Figure~\ref{fig:architecture} shows the system considered here. The uninterruptible power supply (UPS) supports continuity of electricity supply to the IT load. A chiller supplies the computer room air conditioning (CRAC) system and can also charge thermal energy storage (TES), allowing cooling production and cooling delivery to occur at different times. Thermal mass in the equipment and air adds a further dependence on previous operating decisions.

\begin{figure}[t]
\centering\includegraphics[width=\columnwidth]{Figure_1.png}
\caption{Data-centre architecture: IT demand, UPS storage and the coupled chiller, CRAC and TES system.}
\label{fig:architecture}
\end{figure}

These resources interact through both energy flows and operating constraints. Deferring computation reduces IT electricity demand and changes the heat that cooling must remove. TES discharge can reduce direct chiller electricity, while charging the store competes with immediate cooling for chiller capacity. UPS discharge can support IT execution without drawing the same power from the grid. Consequently, one resource can compensate while another completes deferred work or restores its state. Assessing assets independently can miss these compensating actions, or assign the same limited headroom to more than one service.

\subsection{Workloads and the questions addressed}
Computational tasks differ in their tolerance of delay. Interactive services generally require prompt execution, whereas eligible batch tasks can be scheduled within agreed windows~\cite{Cao2022ApEn,Liu2024Access}. This distinction depends on the application and its service-level agreement (SLA), rather than being a fixed property of all workloads of a given type. The model therefore treats workload flexibility as an input: an assumed share of arriving work can be rescheduled within specified limits, and the remainder executes at arrival. It preserves computational work and the modelled timing constraints; it does not represent all aspects of application performance.

Two questions guide the analysis. First, what combinations of power deviation and duration can the integrated facility deliver, and which resources supply them? Second, what electricity-cost savings arise from exploiting its operational flexibility over a full year, and how do those savings change with resource assumptions? The first question concerns physical capability. The second concerns economic operation under the modelled price signal. The annual schedule connects the two by providing the operating state from which additional grid deviations are assessed.

The paper makes three contributions. It presents an integrated representation of workload timing, UPS operation and cooling/TES that retains the couplings relevant to facility-level flexibility. It uses this representation to characterise state-dependent power deviations with post-event recovery, and to explain their internal composition. Finally, it evaluates annual electricity-cost savings and resource sensitivities, using the UPS results to identify the additional investment that the incremental savings could support. The annual rolling procedure supplies the temporal coverage needed for this valuation; it is not proposed as a new optimisation algorithm.

Section~\ref{sec:literature} places these contributions in the literature. Section~\ref{sec:methodology} describes the integrated model and assessment, and Section~\ref{sec:case-study} defines the case study. Section~\ref{sec:results} presents the flexibility characteristics, annual economic results and design implications. Section~\ref{sec:conclusion} concludes.
